\chapter{Infrastructure as an Analytical Problem}
Publication infrastructures are often most visible at moments of breakdown. This synthetic chapter demonstrates a research-monograph hierarchy with sustained prose, notes, and numbered sections.\footnote{The note is synthetic and exists only to demonstrate scholarly page rhythm.}

\section{From artefact to system}
Treating a publication as an isolated artefact obscures the standards, classifications, software, labour, and institutional routines that make it possible.

\subsection{Visibility and maintenance}
Maintenance work rarely appears in the final object. An infrastructural perspective therefore asks not only what a system produces, but how stability is achieved and where responsibility is located.

\section{A working proposition}
The analytical advantage of this approach is that technical choices can be studied as organised relations rather than as neutral background conditions.

\PSEIndex{infrastructure}
